\section{Conclusiones y líneas futuras}
\label{cap5:conclusiones}

Este trabajo ha plasmado la investigación realizada en torno a las redes neuronales compuestas por capas convolucionales y recurrentes aplicadas a la predicción del tráfico en red con \textit{online learning}. Se han analizado varias arquitecturas de redes neuronales con una capa de neuronas LSTM precedida por una cantidad variable entre ninguna y tres capas CNN. Se ha realizado un extenso proceso de validación basado en la estrategia \textit{grid search} sobre el dataset Milán y se ha obtenido el modelo con mayor precisión utilizando el NRMSE como métrica de rendimiento.

Como regla general, los modelos con \textit{online learning} son ampliamente superiores a aquellos sin \textit{online learning}. Los cambios en el patrón de los datos pasan completamente desapercibidos para los modelos sin \textit{online learning}. En cambio, los modelos con \textit{online learning} presentan una capacidad de adaptación extraordinaria.

La implementación de capas convolucionales al modelo de predicción de series temporales basado únicamente en capas de neuronas LSTM ha dado resultado, reduciendo los errores cometidos en las predicciones. Respecto al modelo original cedido por el Dr. D. Juan Pablo Casaseca de la Higuera, se ha conseguido reducir el error de predicción desde NRMSE=0,1778 hasta NRMSE=0,142.

Por último, se ha alcanzado la optimización del modelo original a través de funciones altamente eficientes proporcionadas por la librería \texttt{Tensorflow}. Gracias a ello se ha conseguido reducir considerablemente el tiempo de ejecución requerido por los modelos para los procesos de \textit{backpropagation} y \textit{online learning}.
%%%%OBJETIVOS DEL TFG%%%%%%%%%%%%

Para próximos trabajos se plantea la posibilidad de implementar los modelos descritos en este trabajo como redes de grafos orientadas a la predicción de series temporales, más concretamente a la predicción de tráfico de red.